\honest{Your Strength as a Rationalist}{Eliezer Yudkowsky}{2007}

Long ago, in a chatroom, someone said that a friend of his had chest pains and called an ambulance, and the paramedics told the friend it was nothing and left. Now the pains were getting worse. What should the friend do? I warn you that my memory of this may be imprecise.

The story confused me. I remembered reading that paramedics in New York must take homeless people to hospital even on their 27th call, and that emergency rooms are legally obliged to treat anyone. So anyone reporting chest pains should have been ``hauled off by an ambulance instantly.'' I give no source for any of this. \nb{The conclusion is wrong. A competent adult can refuse to be taken to hospital, and many do.}

I also remembered that my own doctor never panicked about my symptoms, and was always right; once, when I had chest pains, the doctor explained that I was describing muscle pain, not a heart attack. So I told the stranger that if the paramedics said it was nothing, ``it must really be nothing.'' I had explained the story within my model, though the fit ``still felt a little forced.'' His friend's chest pains were getting worse. \nb{If the story had been true, this advice could have delayed treatment for a heart attack. The post does not mention this.}

Later the stranger came back and said his friend had made the whole thing up. So I conclude that my confusion was a clue that the story was false, and that I threw the clue away. \nb{The only evidence that the story was false is the word of the same stranger who told it.}

Perhaps I should have realized that a friend of a stranger in a chatroom is less reliable than a published journal article. But belief is easier than disbelief: we believe by instinct, and disbelief takes effort. I had all the information I needed, I even noticed the problem, and I ignored it.

Then the principle: a model is useful for what it cannot explain, and a model that forbids nothing tells you nothing. \nb{This is true, and it is Karl Popper's idea.}

Then the thesis: ``Your strength as a rationalist is your ability to be more confused by fiction than by reality.'' If you can explain any outcome equally well, you know nothing. \nb{A true report that clashes with a wrong model is just as confusing as a lie, and the feeling does not say which it is.}

That feeling of a forced fit is one of the most important a truthseeker can have. I wish it came as a glowing neon sign reading ``Either Your Model Is False Or This Story Is Wrong.'' \nb{In this case both were true: the story was made up, and the model of ambulances was wrong.}

Two footnotes. The first claims that hospitals are closing their emergency rooms and that economists are being ignored. I give no source. \nb{The footnote has nothing to do with the argument.} In the second I cite evidence that we believe what we hear by default: a post by McCluskey on the group blog \textsc{Overcoming Bias}, and a 1993 paper by Gilbert and colleagues.
